\section{Domain Overview}
\label{sec_domain_overview}

Autonomous penetration testing -- the use of Large Language Model (LLM) agents to independently reconnoitre, identify, and exploit weaknesses in a target system without step-by-step human direction -- has emerged as an active research area only in the last two to three years. The domain sits at the intersection of two older fields: Vulnerability Assessment and Penetration Testing (VAPT), which traditionally relies on human operators following structured methodologies (reconnaissance, weakness identification, exploitation, and reporting), and LLM agent research, which studies how language models can plan, use tools, and act over multiple steps toward a goal. RedGrid is situated in this intersection: it is concerned with how an LLM-driven agent, or a coordinated set of agents, can carry out the VAPT workflow autonomously against realistic, benchmarked targets.

A useful way to frame the domain is around the two capabilities an autonomous VAPT agent must exhibit simultaneously. The first is \emph{exploration}: the agent must search a large space of candidate weaknesses across an unfamiliar attack surface without being told in advance which one is exploitable. The second is \emph{reasoning about prerequisites}: many real exploits are not single actions but chains, where one weakness (for example, an information disclosure or a misconfigured endpoint) must be established before a second, more consequential weakness becomes reachable. Recent benchmarking work suggests that these two capabilities are evaluated very differently in the literature, and that most existing systems have been built to strengthen one at the expense of the other -- a pattern that motivates a closer reading of the existing systems and benchmarks surveyed in this chapter.

The chapter that follows first surveys the systems that have been proposed for this problem (Section~\ref{sec_existing_systems}), then draws out the recurring architectural choices across them (Section~\ref{sec_comparative_analysis}), before identifying, at a preliminary level, where the surveyed literature appears to leave a gap (Section~\ref{sec_research_gap}). As this is an early-stage report, the discussion below reflects the literature review conducted so far rather than a finalized research direction; the survey will continue to be refined over the course of the thesis.
